\section{Introduction}\label{sec:intro}

Training a neural network on a digital processor moves every weight from memory to arithmetic and back at every update.
In 1960, Widrow and Hoff took a different route: they stored each weight of their adaptive ADALINE neuron in a
\emph{memistor}, an electrochemical cell whose conductance is set by the time integral of a current through a third
terminal~\cite{widrow1960adaline,widrow1960switching,widrow1990}. The weight was read without being disturbed, and it was
adapted in place. The device was macroscopic, the cells were wet, and no multilayer learning rule existed yet, so the line ended
with the Madaline hardware of the early 1960s~\cite{widrow1990}. That machine was a working analog learning system.

Sixty years later, three-terminal electrochemical devices have returned at the nanoscale. Lithium-ion
ECRAM~\cite{tang2018,fuller2017}, CMOS-compatible metal-oxide ECRAM~\cite{kim2019}, protonic programmable
resistors~\cite{onen2022,xu2026} and open-loop programmable arrays~\cite{chen2023} all decouple the read path from the write path,
just as the memistor did. They report hundreds to thousands of conductance levels, nanosecond pulses and femtojoule-scale write
energies. In parallel, algorithms for training on imperfect analog weights have matured: the stochastic pulse-coincidence update
of the resistive processing unit (RPU)~\cite{gokmen2016}, Tiki-Taka~\cite{gokmen2020} and its successors
TTv2~\cite{gokmen2021}, chopped TTv2 and AGAD~\cite{rasch2024}, and residual learning with convergence guarantees~\cite{wu2025}.

That electrochemical weights can train small networks is already established: Chen et al.\ trained in situ on an ECRAM array and
simulated larger networks~\cite{chen2023}, and Rasch et al.\ studied reference correction, retention and device requirements for
the TT family~\cite{rasch2024}. Most recently, Wu et al.~\cite{wu2026scale} simulated ECRAM-calibrated Transformers of up to about
123 million parameters with finite converters and device nonidealities, keeping the weight-gradient path digital. Byun et al.~\cite{byun2024} derived device specifications for TTv1 and TTv2 by
simulation; relative to that analysis, this study covers the chopped algorithms c-TTv2 and AGAD and a catalog of published
electrochemical devices, and reports seed-paired intervals against absolute margins, a language model beside the image
tasks and a test of the aggregate update against explicit pulse streams; its own sensitivity analysis is narrower. Our
contribution is narrower and specific: \emph{a reproducible, seed-matched simulation comparison} of plain SGD and three
device-tolerant algorithms on a catalog of published (or, where nothing is published, labelled hypothetical) electrochemical
devices, with the weight update done entirely in the devices, on three small tasks; a sensitivity analysis against absolute
accuracy lines; and a \emph{conditional} energy projection. Results are stated against two digital baselines: SGD with exact
updates at the same learning rate, and a separately tuned Adam.

What this paper provides:
\begin{enumerate}
\item \textbf{A simulator with stated limits.} Reads are ideal (fp32); only the update is analog (Fig.~\ref{fig:arch}). It reproduces selected
      published outcomes of Gokmen and Haensch 2020 under this simplified read model, and the offset-robustness ordering of
      Rasch et al.\ 2024 (Section~\ref{sec:validation}). After an independent review it was repaired (separate random streams, a
      reset made of pulses, the AIHWKIT device presets mapped one to one) and each repair is tested (Section~\ref{sec:methods},
      Appendix~\ref{app:revision}).
\item \textbf{A seed-matched comparison} of SGD, TTv2, chopped TTv2 and AGAD on five device specifications, three tasks (MNIST, a
      CIFAR-10 subset, a byte-level language model on human-written text) and five seeds, with every arm of a seed sharing its
      initial weights and data order (Section~\ref{sec:main}).
\item \textbf{A width study} with per-arm learning rates and an equivalence test (Section~\ref{sec:scale}).
\item \textbf{An exploratory device sensitivity analysis} against absolute accuracy lines, one property at a time with three
      seeds (not fabrication tolerances), with the asymmetry sweep repeated at a fixed symmetry point
      so that asymmetry is not confounded with lost weight range (Section~\ref{sec:spec}).
\item \textbf{A conditional energy projection} for the weight-matrix operations, reported only where the device reaches the
      stated quality (Section~\ref{sec:energy}).
\item \textbf{A proposed roadmap} of device experiments and a proposed, unbuilt bench replication protocol of the 1960
      memistor (Section~\ref{sec:roadmap}, Appendix~\ref{app:bench}).
\end{enumerate}

Every experiment was planned in a committed plan before it ran; the plans are internal, timestamped by the repository, not an
external registry. The original runs and the revision runs are kept apart, and every deviation is recorded
(Section~\ref{sec:availability}). All results are simulations, and every device parameter carries a provenance label
(measured, fitted, claimed, estimated or assumed).
